\section{Related Work}
\label{sec:related}

\para{Selection bias in multiple-choice answering.}
\citet{zheng2024selectors} show that LLMs have a selection bias in multiple-choice questions, attribute much of it to a preference for particular label \emph{tokens} rather than positions, and propose cyclic permutation as a cheap way to estimate and remove it.
\citet{pezeshkpour2024order} report large accuracy swings when options are reordered.
Building on these studies, we use cyclic rotation as a measurement tool rather than a debiasing method, and we ask a narrower question: how many accuracy points does position add to a score under explicit definitions of inflation?
Unlike comparisons across benchmarks with different numbers of options, we vary the option count on the same questions, so format is not confounded with content.

\para{Reading out answers.}
\citet{robinson2023mcqa} study prompting LLMs with all options and reading the probability of the answer symbol.
\citet{wang2024myanswer} show that first-token probabilities can disagree with the text an instruction-tuned model generates.
We use a first-token read-out because it is deterministic and gives a full distribution over labels, and we check it against greedy generation on a subsample (\secref{sec:readout}).
The check also identifies the one model and the one label set for which the read-out is unreliable.

\para{Benchmarks and uncertainty in evaluation.}
We use \mmlu~\citep{hendrycks2021mmlu}, \csqa~\citep{talmor2019commonsenseqa}, and \boolq~\citep{clark2019boolq} because they cover 4-option, 5-option, and yes/no formats.
\citet{miller2024errorbars} argues for reporting uncertainty on evaluation scores and for clustering by question when responses are not independent.
We follow this practice: every interval we report is a bootstrap over questions that resamples all rotations of a question together.

\para{Difficulty and uncertainty.}
\citet{zheng2024selectors} relate selection bias to model uncertainty.
We revisit this link with a stratification that does not reuse the responses being analyzed: we define the difficulty of a question for one model from the other six models.
